\documentclass[10pt,twocolumn,letterpaper]{article}

% CVPR packages
\usepackage[final]{cvpr}
\usepackage{times}
\usepackage{epsfig}
\usepackage{graphicx}
\usepackage{amsmath}
\usepackage{amssymb}
\usepackage{booktabs}
\usepackage{color}
\usepackage{url}

% Include other packages as needed
\usepackage{caption}
\usepackage{subcaption}
\usepackage{algorithm}
\usepackage{algpseudocode}

% If you comment hyperref and then uncomment it, you should delete
% egpaper.aux before re-running latex.  (Or just hit 'q' on the first latex
% run, let it finish, and you should be clear).
\usepackage[pagebackref,breaklinks,colorlinks]{hyperref}

% Support for easy cross-referencing
\usepackage[capitalize]{cleveref}
\crefname{section}{Sec.}{Secs.}
\Crefname{section}{Section}{Sections}
\Crefname{table}{Table}{Tables}
\crefname{table}{Tab.}{Tabs.}

% Page style
\cvprfinalcopy % *** Uncomment this line for the final submission

\def\cvprPaperID{*****} % *** Enter the CVPR Paper ID here
\def\confName{CVPR}
\def\confYear{2025}

\begin{document}

% Title
\title{Paper Title}

% Authors
\author{
Anonymous Authors\\
Institution\\
\tt\small anonymous@institution.edu
}

\maketitle
% Remove page # from the first page of camera-ready.
\thispagestyle{empty}

%%%%%%%%% ABSTRACT
\begin{abstract}
Your abstract here. Summarize the problem, proposed method, and key results in 100-150 words.
\end{abstract}

%%%%%%%%% BODY TEXT
\section{Introduction}
\label{sec:intro}

Introduction here. Motivate the problem, state your contributions clearly, and provide an overview of the paper structure.

\textbf{Contributions:}
\begin{itemize}
    \item First contribution
    \item Second contribution
    \item Third contribution
\end{itemize}

\section{Related Work}
\label{sec:related}

Review related work here. Position your method relative to existing approaches.

\section{Method}
\label{sec:method}

Describe your method in detail. Include:
\begin{itemize}
    \item Problem formulation
    \item Technical details with equations
    \item Architecture diagrams (use figures)
    \item Algorithm pseudocode if applicable
\end{itemize}

\subsection{Overview}

Provide an overview of your approach.

\subsection{Detailed Design}

Describe the technical details.

\begin{equation}
    \mathcal{L} = \mathcal{L}_{cls} + \lambda \mathcal{L}_{reg}
\end{equation}

\section{Experiments}
\label{sec:experiments}

\subsection{Experimental Setup}

Describe:
\begin{itemize}
    \item Datasets used
    \item Evaluation metrics
    \item Implementation details (optimizer, learning rate, etc.)
    \item Hardware used
\end{itemize}

\subsection{Main Results}

Present your main results with tables and figures. Compare against baselines and state-of-the-art methods.

\begin{table}[t]
\centering
\caption{Comparison with state-of-the-art methods on CIFAR-100.}
\label{tab:main_results}
\begin{tabular}{lccc}
\toprule
Method & Params & Top-1 Acc & Top-5 Acc \\
\midrule
Baseline & 11M & 75.0 & 92.5 \\
Proposed & 10M & \textbf{78.5} & \textbf{94.2} \\
\bottomrule
\end{tabular}
\end{table}

\subsection{Ablation Studies}

Conduct ablation studies to validate design choices.

\subsection{Analysis}

Provide additional analysis, visualizations, and insights.

\section{Conclusion}
\label{sec:conclusion}

Summarize your work, restate key findings, and discuss limitations and future directions.

\textbf{Limitations:} Discuss limitations of your method.

\textbf{Future Work:} Suggest directions for future research.

%%%%%%%%% REFERENCES
{\small
\bibliographystyle{ieee_fullname}
\bibliography{references}
}

\end{document}
